\input{../tex-inputs/latex-leseansicht-vorspann}

\section[Arthur und Olga Schnitzler an Gustav Schwarzkopf, 4. 9. 1908]{L04382 Arthur und Olga Schnitzler an Gustav Schwarzkopf, 4. 9. 1908}
\nopagebreak\mylabel{L04382v}
\rehead{ }\normalsize\beginnumbering\briefempfaengerindex{Schwarzkopf, Gustav@\textsc{Schwarzkopf, Gustav}!zzzSchnitzler, Olga@\emph{von Olga Schnitzler}!1908-09-041@{4. 9. 1908}|(be}\briefempfaengerindex{Schwarzkopf, Gustav@\textsc{Schwarzkopf, Gustav}!zzzSchnitzler, Arthur@\emph{von Arthur Schnitzler}!1908-09-041@{4. 9. 1908}|(be}
\toendnotes[C]{\smallbreak\pagebreak[2]}
\correspDesc{Versand  durch Arthur Schnitzler, Olga Schnitzler am 4. 9. 1908 in San Martino di Castrozza
\newline{}Übermittlung  am 5. 9. 1908 in San Martino di Castrozza
\newline{}Erhalt  durch Gustav Schwarzkopf im Zeitraum [5. 9. 1908
                  – 9. 9. 1908?] in Wien}\toendnotes[C]{\smallbreak}
\Standort{DLA, A:Schnitzler, HS.1985.1.1897.}
\physDesc{Bildpostkarte, 342 Zeichen
\newline{}Handschrift Arthur Schnitzler: Bleistift, deutsche Kurrent
\newline{}Handschrift Olga Schnitzler: Bleistift, lateinische Kurrent
\newline{}Versand: Stempel: »\nobreak{}\oindex{San Martino di Castrozza@\textbf{San Martino di Castrozza}|pwk}San Martino di Castrozza, \textcolor{gray}{5}/9 \textcolor{gray}{08}\nobreak{}«.  }\toendnotes[C]{\smallbreak}\pstart{}{\pb}Hrn \textsc{Gustav Schwarzkopf}\pend{}\pstart{}\textsc{Wien I}\oindex{I., Innere Stadt@\textbf{I., Innere Stadt}|pw}\pend{}\pstart{}\textsc{Tiefer Graben 23}\oindex{Wien@\textbf{Wien}!I., Innere Stadt@\textbf{I., Innere Stadt}!Tiefer Graben 23@\textbf{Tiefer Graben 23}|pw}\pend{}{\bigskip}
\pstart
           \noindent{}\centering{}{\pb}\textcolor{gray}{\textbf{S. Martino di Castrozza\oindex{San Martino di Castrozza@\textbf{San Martino di Castrozza}|pw} (1465 m) mit Cavalazza\oindex{Cavallazza@\textbf{Cavallazza}|pw}.}}\pend
           \vspace{1em}
\pstart
           \noindent{}{\pb}lieber Guſtav,{ }\label{K_L04382-1v}\edtext{bis vorgeſtern}{\lemma{\textnormal{\emph{bis vorgestern}}}\Cendnote{\textnormal{Siehe A. S.: \emph{Wiener Schnitzler}, 2. 9. 1908.}}}\label{K_L04382-1} waren wir
               in Seis\oindex{Seis am Schlern@\textbf{Seis am Schlern}|pw} u{ }ſind nun auf allmäliger Rückreiſe, die
               auch noch \label{K_L04382-2v}\edtext{über München\oindex{München@\textbf{München}|pw}}{\lemma{\textnormal{\emph{über München}}}\Cendnote{\textnormal{Siehe A. S.: \emph{Tagebuch}, 7. 9. 1908.}}}\label{K_L04382-2} führen{ }ſoll. Der So{\geminationm}er war ſehr ſchön –  wo ſind Sie geweſen?
               Auf baldgs Wiederſehn – herzliche{\\}Grüße! \spacefill\mbox{Ihr A.}\pend
           \selectlanguage{ngerman}\vspace{1em}
\pstart
           \noindent{}{[}hs. O. Schnitzler:{]} \label{K_L04382-3v}\edtext{Lisl\pwindex{Steinrück, Elisabeth 19.\,11.\,1885 – 7.\,4.\,1920 Partenkirchen@\textsc{Steinrück, Elisabeth} (19.\,11.\,1885 – 7.\,4.\,1920 Partenkirchen)|pw}{ }heiratet\eventindex{Hochzeit von Elisabeth Gussmann und Albert Steinrück, September 1908@Hochzeit von Elisabeth Gussmann und Albert Steinrück, September 1908|pwv}}{\lemma{\textnormal{\emph{Lisl heiratet}}}\Cendnote{\textnormal{Elisabeth
                     Gussmann\pwindex{Steinrück, Elisabeth 19.\,11.\,1885 – 7.\,4.\,1920 Partenkirchen@\textsc{Steinrück, Elisabeth} (19.\,11.\,1885 – 7.\,4.\,1920 Partenkirchen)|pwk}{ }heiratete\eventindex{Hochzeit von Elisabeth Gussmann und Albert Steinrück, September 1908@Hochzeit von Elisabeth Gussmann und Albert Steinrück, September 1908|pwkv} im September 1908{ }Albert Steinrück\pwindex{Steinrück, Albert 20.\,5.\,1872 Wetterburg – 11.\,2.\,1929 Berlin@\textsc{Steinrück, Albert} (20.\,5.\,1872 Wetterburg – 11.\,2.\,1929 Berlin)|pwk}.}}}\label{K_L04382-3}
               in d. nächsten Tagen.\pend
           
\pstart
           {\pb}Herzliche \uline{Grüsse} und auf
               Wiedersehen! \spacefill\mbox{OlgaS.}\pend
           
\pstart
           4. Sept. 08.\pend
           \selectlanguage{ngerman}\endnumbering\briefempfaengerindex{Schwarzkopf, Gustav@\textsc{Schwarzkopf, Gustav}!zzzSchnitzler, Olga@\emph{von Olga Schnitzler}!1908-09-041@{4. 9. 1908}|)be}\briefempfaengerindex{Schwarzkopf, Gustav@\textsc{Schwarzkopf, Gustav}!zzzSchnitzler, Arthur@\emph{von Arthur Schnitzler}!1908-09-041@{4. 9. 1908}|)be}\mylabel{L04382h}
\begin{anhang}
\end{anhang}\newcommand{\dateiname}{L04382}\newcommand{\titel}{Arthur und Olga Schnitzler an Gustav Schwarzkopf, 4. 9. 1908}\newcommand{\editorInnen}{Herausgegeben von Jahnke, SelmaMüller, Martin Anton}\input{../tex-inputs/latex-leseansicht-abspann}
